\section{The construction and the rank obstruction}\label{sec:construction}

\subsection{The geometric criterion}

A \emph{finite Euclidean triangle complex} will mean a finite graph with finitely many Euclidean triangles attached side-to-edge along their full boundaries, with positive compatible lengths and isometric gluings. Each closed edge and triangle is required to embed; different cells may intersect in a union of faces. The angular link at a vertex has one vertex for each incident edge germ and one edge for each triangle corner, of length equal to that corner's angle. A walk is \emph{reduced} if it has no immediate reversal, and \emph{cyclically reduced} if this also holds across its closing seam. Link distances between different components are infinite.

\begin{theorem}[Triangle complex criterion]\label{thm:triangle-complex}
Let $K$ be a connected finite Euclidean triangle complex. Suppose there is $\eta>0$ such that every nonempty cyclically reduced closed edge walk in each vertex link has angular length at least $2\pi+\eta$. Then $\pi_1(K)$ is torsion-free and word-hyperbolic.

Moreover, a nonempty closed edge path in $K$ is not nullhomotopic if every cyclic turn has link distance at least $\pi$. Here a turn compares the direction back along the arriving edge with the direction of the departing edge.
\end{theorem}

The criterion uses the classical angular-link approach to curvature;
see \cite[Theorem~II.5.2 and Lemma~II.5.6]{BH1999} for its general metric framework. We give
the complete proof of the precise statements above in
Sections~\ref{sec:geometry-foundations} and~\ref{sec:geometry-thinness},
including nontriviality for the specified turns in the original
Euclidean metric. We first use the criterion to specify the group.

\subsection{Punctured lines with few short cycles}

We need many line labels whose punctured incidence graph has large girth,
while almost every point still meets enough labels to admit the two-block
partition below. We use the probabilistic alteration method: random
sampling is followed by deletion of labels on short cycles, as in the
high-girth argument of Erd\H{o}s \cite[pp.~35--37]{Erdos1959}. The marked
line data and the exceptional-set estimates are constructed below.

Fix
\begin{equation}\label{eq:fixed-parameters}
 h=20,\qquad c=100^{-h},\qquad r\in\N\quad\text{with }cr\ge100.
\end{equation}
All limits below are as the prime $q$ tends to infinity with these constants fixed. Put $P=\F_q^h$ and $m=q^h$, and choose a subset $S\subset\F_q$ of size $\lfloor q/100\rfloor$.

\begin{lemma}\label{lem:incidence}
For all sufficiently large primes $q$, there are $N=(c+o(1))m$ affine lines $\ell_i\subset P$, indexed by a set $I$, with distinct marked points $a_i\in S^h\cap\ell_i$, and a set $E\subset P$ such that:
\begin{enumerate}[label=\textup{(\roman*)}]
\item the bipartite graph on $P\sqcup I$, with edges $z\sim i$ when $z\in\ell_i\setminus\{a_i\}$, has no cycle of length at most $12$;
\item $|I(z)|\ge4r^2$ for $z\notin E$, where
\[
 I(z)=\{i\in I:z\in\ell_i\setminus\{a_i\}\};
\]
\item $|E|=o(m)$.
\end{enumerate}
In particular we can fix such data with
\begin{equation}\label{eq:strict-budget}
 q>r,\qquad \frac{Nr}{4}>2m+r|E|.
\end{equation}
\end{lemma}

The second inequality in~\eqref{eq:strict-budget} is the numerical gap
that the final rank comparison will contradict.

\begin{proof}
Initially use every point of $S^h$ as a mark, independently choosing a uniform direction for the line through it. The number of directions is
\[
 J=\frac{q^h-1}{q-1}.
\]
For a fixed $z\in P$, each mark other than $z$ gives an incidence with probability $1/J$. Its degree is therefore binomial, with mean $\mu_z\sim cq$ uniformly in $z$ and variance at most $\mu_z$. For large $q$, $\mu_z\ge8r^2$, so the second-moment inequality gives
\[
 \Pr\bigl(|I(z)|<4r^2\bigr)
 \le \frac{\mu_z}{(\mu_z-4r^2)^2}=O(q^{-1}).
\]
The expected number of low-degree points is $O(m/q)$. By Markov's inequality, with probability tending to one there are at most $m q^{-1/2}$ such points.

For $2\le s\le6$, count simple cycles with $s$ point vertices and $s$ line labels. There are at most $m^s$ ordered choices of the points. Consecutive points force a geometric line, on which there are at most $q$ possible marks. For distinct line labels, the prescribed directions occur with probability $J^{-s}$ by independence. Choices whose mark is a forbidden punctured endpoint can only increase this upper bound. The expected number of these cycles is consequently at most
\[
 m^s q^s J^{-s}\le q^{2s},
\]
since $J\ge q^{h-1}$. This argument also allows different labels to specify the same geometric line: their random directions remain independent. Summing over $s$ shows that, with probability tending to one, the total number of cycles of length at most $12$ is at most $q^{13}$.

Choose an outcome satisfying both bounds. Delete every label appearing in one of those cycles, at most $6q^{13}$ labels, and let $E$ contain the old low-degree points and all points of the discarded lines. Deleting labels creates no cycle, and outside $E$ it changes no degree. Moreover,
\[
 |E|\le m q^{-1/2}+6q^{14}=o(m),\qquad
 N=|S|^h-O(q^{13})=(c+o(1))m.
\]
As $cr\ge100$, these estimates give~\eqref{eq:strict-budget} for large enough $q$: eventually $N/m\ge c/2$ and $r|E|/m\le1$, whereas $cr/8\ge12.5>3$.
\end{proof}

Fix the finite data in Lemma~\ref{lem:incidence} once and for all. No subsequent choice of a finite quotient will change $q$ or any of these data.

\subsection{Blocks and triangular relations}

For $z\notin E$, partition $I(z)$ into two blocks $I_b$ whose sizes factor as
\[
 |I_b|=s_b t_b,\qquad s_b,t_b\ge r.
\]
Indeed, for $d=|I(z)|\ge4r^2$, choose $n_2\in[r,2r-1]$ with $n_2\equiv d\pmod r$. Then
\[
 d=r n_1+(r+1)n_2,\qquad
 n_1=\frac{d-(r+1)n_2}{r}\ge\frac{2r^2-r+1}{r}\ge r.
\]
Use blocks of sizes $r n_1$ and $(r+1)n_2$. Index blocks at different points separately and write $p(b)$ for the point of block $b$. At points of $E$ make no blocks. For each block fix a bijection
\begin{equation}\label{eq:block-coordinates}
 i\longmapsto\bigl(\alpha_b(i),\beta_b(i)\bigr)
 \in\Z/s_b\Z\times\Z/t_b\Z.
\end{equation}

Let $[r]=\{0,\ldots,r-1\}$. Form a complex $K$ with vertices $O,V,W$ and edges
\[
\begin{array}{ll}
 x_{iuv}:V\longrightarrow W & (i\in I,\ u,v\in[r]),\\[2pt]
 L_{b,u,j}:O\longrightarrow V & (u\in[r],\ j\in\Z/s_b\Z),\\[2pt]
 R_{b,v,k}:O\longrightarrow W & (v\in[r],\ k\in\Z/t_b\Z).
\end{array}
\]
For every $b$, $i\in I_b$, and $u,v\in[r]$, attach a triangle with sides
\begin{equation}\label{eq:triangle-sides}
 x_{iuv},\qquad L_{b,u,\alpha_b(i)+v},\qquad
 R_{b,v,\beta_b(i)+u}.
\end{equation}
Give each triangle angles $3\pi/5,\pi/5,\pi/5$ at $O,V,W$, respectively, using one fixed shape. All edge lengths are compatible; assign the same appropriate lengths to unused edges as well. Closed cells embed because the three vertices of every triangle are distinct. There are regular points and hence blocks, so $K$ is connected.

Figure~\ref{fig:construction-triangle} shows one of these triangles.
The larger angle at $O$ will control the four-edge cycles in its link;
the smaller angles at $V$ and $W$ will be paired with the large girth of
the incidence graph.

\begin{figure}[htbp]
\centering
\begin{tikzpicture}[scale=1.1, every node/.style={font=\small}]
  \coordinate (O) at (0,1.6);
  \coordinate (V) at (-2.20221,0);
  \coordinate (W) at (2.20221,0);
  \fill[gray!8] (O)--(V)--(W)--cycle;
  \draw[-{Stealth[length=2mm]}] (O)--node[above left] {$L$} (V);
  \draw[-{Stealth[length=2mm]}] (O)--node[above right] {$R$} (W);
  \draw[-{Stealth[length=2mm]}] (V)--node[below] {$x$} (W);
  \draw (O) ++(216:.43) arc[start angle=216,end angle=324,radius=.43];
  \draw (V) ++(0:.62) arc[start angle=0,end angle=36,radius=.62];
  \draw (W) ++(144:.62) arc[start angle=144,end angle=180,radius=.62];
  \node at (0,.92) {$3\pi/5$};
  \node at (-1.33,.28) {$\pi/5$};
  \node at (1.33,.28) {$\pi/5$};
  \node[above] at (O) {$O$};
  \node[left] at (V) {$V$};
  \node[right] at (W) {$W$};
\end{tikzpicture}
\caption{One attached triangle, with $L,R,x$ abbreviating the indexed
edges in \eqref{eq:triangle-sides}. The oriented boundary is
$LxR^{-1}$. Many distinct edges and triangles share the same three
vertices in $K$; the picture shows only one triangle.}
\label{fig:construction-triangle}
\end{figure}

\begin{lemma}\label{lem:links}
The complex $K$ satisfies Theorem~\ref{thm:triangle-complex} with $\eta=2\pi/5$. At either $V$ or $W$, the link distance between distinct $x$-directions with the same label $i$ is at least $14\pi/5$, or is infinite.
\end{lemma}

\begin{proof}
At $O$, each block gives a complete bipartite link on its $L$- and $R$-directions. For fixed $u,v,j,k$, the equations
\[
 \alpha_b(i)=j-v,\qquad \beta_b(i)=k-u
\]
specify exactly one $i\in I_b$. Thus there is one link edge per pair of directions, and every nonempty cyclically reduced closed walk has at least four edges. Its angular length is at least $12\pi/5$.

Project the link at $V$ to the incidence graph of Lemma~\ref{lem:incidence} by sending an $x_{iuv}$-direction to $i$ and an $L_{b,u,j}$-direction to $p(b)$. This map is locally injective on edges. At $x_{iuv}$ there is at most one incident triangle over any point $z$, because $i$ lies in at most one block at $z$. At $L_{b,u,j}$, fixing $i$ determines at most one $v\in[r]$ from $j=\alpha_b(i)+v$, since $s_b\ge r$. The link at $W$ has the same property using $t_b\ge r$ and the $R$-directions. Every cyclically reduced link walk therefore has a reduced image and at least $14$ edges, of angular length at least $14\pi/5$.

A reduced path between distinct $x$-directions with the same $i$ projects to a positive closed walk reduced at every internal vertex. It need not be reduced at its closing seam, but it still contains a simple cycle: choose a shortest positive subwalk with equal endpoints. A subwalk of length two would be an internal reversal, and one of length one is impossible in this graph. The path therefore has at least $14$ edges. A shortest link path is reduced, proving the distance assertion.
\end{proof}

Set $G=\pi_1(K,O)$. Theorem~\ref{thm:triangle-complex} and Lemma~\ref{lem:links} show that $G$ is torsion-free and word-hyperbolic. It remains to exclude residual finiteness.

Choose a path from $O$ to each vertex and associate to an oriented edge $e$ the based loop given by the path to its initial vertex, then $e$, then the reverse of the path to its terminal vertex. Write its element as $g_e$, using multiplication in path order, and abbreviate $g_{iuv}=g_{x_{iuv}}$. The triangles give
\begin{equation}\label{eq:group-factorization}
 g_{iuv}=g_{L_{b,u,\alpha_b(i)+v}}^{-1}
               g_{R_{b,v,\beta_b(i)+u}}\qquad(i\in I_b).
\end{equation}
Furthermore,
\begin{equation}\label{eq:rectangles}
 g_{iuv}g_{iu'v}^{-1}g_{iu'v'}g_{iuv'}^{-1}\ne1
 \qquad(u\ne u',\ v\ne v').
\end{equation}
Indeed, after cancelling the chosen connecting paths this word is conjugate to a closed four-edge path. Each of its four turns compares distinct same-$i$ directions in a link at $V$ or $W$, and so has distance at least $\pi$ by Lemma~\ref{lem:links}. The closed-path assertion of Theorem~\ref{thm:triangle-complex} applies.

Let $\mathcal R\subset G$ be the finite set of rectangular words in
\eqref{eq:rectangles}, with $i\in I$, $u\ne u'$ and $v\ne v'$ in $[r]$.
We have proved $1\notin\mathcal R$. The remaining argument shows that
no finite quotient can keep every member of $\mathcal R$ nonidentity.

\subsection{Rank over the reals and over a fixed prime field}

The following elementary observation makes it possible to choose the characteristic independently of a prospective finite quotient.

\begin{lemma}[Rank under reduction]\label{lem:rank-reduction}
Let $M$ be an integer matrix whose rows have Euclidean norm at most $\sqrt r$, where $r\ge1$. For every prime $q>r$,
\[
 \rank_{\F_q}M\ge\tfrac12\rank_{\R}M.
\]
\end{lemma}

\begin{proof}
Let $d=\rank_{\R}M$; the assertion is immediate for $d=0$. A nonzero $d\times d$ minor $B$ has
\[
 1\le|\det B|\le r^{d/2},
\]
by Hadamard's determinant inequality, obtained by orthogonalizing its rows. If $d'=\rank_{\F_q}B$, choose $d'$ rows spanning its row space modulo $q$. Subtract integer lifts of their linear combinations from the other rows. These determinant-preserving operations make $d-d'$ rows divisible by $q$, whence
\[
 q^{d-d'}\mid\det B,\qquad q^{d-d'}\le r^{d/2}.
\]
Since $q>r$, we obtain $d'\ge d/2$, as required.
\end{proof}

Suppose, for a contradiction, that a finite quotient $Q$ of $G$ keeps every member of $\mathcal R$ nonidentity. Put $D=|Q|$. For an element $g\in G$, let $U(g)$ be its $D\times D$ regular permutation matrix on $Q$, using left multiplication and column vectors. Thus $U$ is a homomorphism and $U(g)^{\mathsf T}=U(g^{-1})$.

For $i\in I$ form the integer matrix of $D\times D$ blocks
\begin{equation}\label{eq:Yi}
 Y_i=\bigl(U(g_{iuv})\bigr)_{u,v\in[r]}.
\end{equation}
The real trace of the regular permutation of a nonidentity element is zero, whereas the identity has trace $D$. Expanding in blocks gives
\begin{align}
 \tr(Y_iY_i^{\mathsf T})&=r^2D,\label{eq:trace-first}\\
 \tr\bigl((Y_iY_i^{\mathsf T})^2\bigr)&=(2r^3-r^2)D.\label{eq:trace-second}
\end{align}
In the second expansion, a term is the trace of
\[
 U\bigl(g_{iuv}g_{iu'v}^{-1}g_{iu'v'}g_{iuv'}^{-1}\bigr).
\]
It equals $D$ when $u=u'$ or $v=v'$, and is zero otherwise by~\eqref{eq:rectangles}. There are $2r^3-r^2$ degenerate index choices.

Cauchy--Schwarz on the positive eigenvalues of $Y_iY_i^{\mathsf T}$ now gives
\[
 \rank_{\R}Y_i\ge
 \frac{r^4D^2}{(2r^3-r^2)D}
 =\frac{r^2D}{2r-1}\ge\frac{rD}{2}.
\]
Each row of $Y_i$ has exactly $r$ unit entries. Lemma~\ref{lem:rank-reduction} therefore yields
\begin{equation}\label{eq:rank-lower}
 \rank_{\F_q}Y_i\ge\frac{rD}{4}.
\end{equation}
There is no restriction on $D$, and in particular no assumption that $q$ does not divide $D$.

\subsection{Local factorizations and the global certificate}

We will express the direct sum of the high-rank matrices $Y_i$ as a sum
of terms with small total rank. The triangle relations give the local
rank bounds; interpolation supplies weights that combine them.

Work henceforth over $\F_q$, and put
\[
 Z(z)=\sum_{i\in I(z)}Y_i\qquad(z\in P).
\]
At a regular point $z\notin E$, each of its two blocks satisfies
\begin{equation}\label{eq:matrix-factorization}
 \sum_{i\in I_b}Y_i
 =\left(\sum_{j\in\Z/s_b\Z}U(g_{L_{b,u,j}}^{-1})\right)_{u\in[r]}
  \left(\sum_{k\in\Z/t_b\Z}U(g_{R_{b,v,k}})\right)_{v\in[r]}.
\end{equation}
The first factor is a block column and the second a block row. For each fixed $u,v$, the shifted coordinates in~\eqref{eq:group-factorization} run through all pairs $(j,k)$ exactly once. The two complete sums are independent of the opposite block index, so~\eqref{eq:matrix-factorization} is a factorization of the whole matrix, with inner dimension $D$. The order of the factors is preserved; no commutativity is used. Hence
\begin{equation}\label{eq:rank-upper}
 \rank Z(z)\le
 \begin{cases}
 2D,&z\notin E,\\
 rD,&z\in E.
 \end{cases}
\end{equation}
For exceptional points we use the ambient matrix dimension. Their incidences have not been deleted from $I(z)$.

We now construct scalar matrix weights whose sum over a punctured line
isolates that line's label. For each $i\in I$, product Lagrange interpolation on $S^h$ supplies a polynomial $f_i$ of total degree at most $h(|S|-1)$ with
\[
 f_i(a_{i'})=\begin{cases}1,&i=i',\\0,&i\ne i'.\end{cases}
\]
Explicitly, for $a_i=(a_{i1},\ldots,a_{ih})$, take
\[
 f_i(X_1,\ldots,X_h)=
 \prod_{t=1}^h\ \prod_{a\in S\setminus\{a_{it}\}}
 \frac{X_t-a}{a_{it}-a}.
\]
Let $f$ be the $N$-entry column of these polynomials and set $H(z)=f(z)f(z)^{\mathsf T}$, an $N\times N$ matrix. The matrices $Z(z)$ have size $rD\times rD$. Then $\rank H(z)\le1$ and $H(a_i)=E_{ii}$, the $i$th diagonal matrix unit.

Every entry of $H$ restricts to a polynomial of degree at most $2h(|S|-1)<q-1$ on an affine line. The finite-field power-sum identity says that such a polynomial sums to zero over $\F_q$. Indeed the constant sum is $q=0$, and for $1\le j<q-1$ there is $a\in\F_q^\times$ with $a^j\ne1$, by the polynomial root bound. Multiplying the sum of $T^j$ by $a^j$ permutes its summands and forces that sum to vanish. Thus
\begin{equation}\label{eq:line-certificate}
 \sum_{z:\,i\in I(z)}H(z)=-H(a_i)=-E_{ii}.
\end{equation}
Only the line's own mark is omitted. Other marks on that line, if any, remain in the sum.

Exchanging finite sums in~\eqref{eq:line-certificate} gives
\begin{equation}\label{eq:tensor-certificate}
 -\bigoplus_{i\in I}Y_i
 =-\sum_{i\in I}E_{ii}\otimes Y_i
 =\sum_{z\in P}H(z)\otimes Z(z).
\end{equation}
The rank of the left side is at least $NrD/4$ by~\eqref{eq:rank-lower}. The image of $H(z)\otimes Z(z)$ lies in the tensor product of the two images, of dimension at most $\rank Z(z)$. Rank subadditivity and~\eqref{eq:rank-upper} therefore give
\[
 \frac{NrD}{4}\le 2mD+r|E|D.
\]
After cancellation of the positive integer $D$, this contradicts~\eqref{eq:strict-budget}.
Thus no finite quotient preserves all of
$\mathcal R$. If each $w\in\mathcal R$ had a nonidentity image in some
finite group, the product of those finitely many maps, followed by taking
its image, would give precisely such a quotient. Consequently there is
one fixed $w\in\mathcal R$ whose image is the identity under every
homomorphism from $G$ to a finite group. Since $w\ne1$ in $G$, the finite
residual of $G$ is nontrivial. This proves Theorem~\ref{thm:main}, subject
only to the geometric criterion proved next.
